
\chapter*{Introduction}

Les données géographiques occupent une place importante dans le fonctionnement
de nombreuses organisations. Elles sont utilisées dans des domaines variés tels
que l'aménagement du territoire, la gestion des réseaux, l'environnement, les
transports ou encore l'urbanisme. Cependant, ces données sont souvent réparties
entre plusieurs bases de données, serveurs, répertoires de fichiers et
applications métiers. Cette dispersion peut rendre leur identification, leur
compréhension et leur réutilisation difficiles.

Isogeo développe des solutions permettant aux organisations de mieux recenser,
documenter et valoriser leur patrimoine de données géographiques. Parmi ces
solutions, le produit Scan permet d'explorer différentes sources de données,
d'identifier les ressources qu'elles contiennent et d'en extraire les
informations techniques nécessaires à leur catalogage. Il peut notamment
analyser des bases PostgreSQL/PostGIS, Oracle et Microsoft SQL Server, des
géodatabases d'entreprise Esri ainsi que plusieurs formats de fichiers
géographiques.

Historiquement, l'ensemble des traitements réalisés par Scan reposait sur des
workbenches FME. Ces traitements permettaient de parcourir les différentes
sources, d'analyser les ressources disponibles et de produire les informations
attendues par le produit. Toutefois, cette organisation rendait le
fonctionnement de Scan dépendant de l'installation et de la licence d'un
logiciel externe.

Dans une démarche d'évolution du produit, Isogeo a donc engagé la migration
progressive de l'ensemble de ces traitements FME vers une solution développée
en Python. L'objectif est de rendre Scan plus autonome, de faciliter la
maintenance de ses traitements et de permettre leur exécution sous la forme
d'un programme Windows ne nécessitant pas l'installation préalable de Python ou
de FME.

Mon alternance s'inscrit dans ce contexte. Ma mission principale consiste à
poursuivre cette migration en analysant les traitements existants, leurs
entrées, leurs sorties et les règles métier qu'ils appliquent, puis en
développant des scripts Python capables de produire des résultats équivalents.
Elle comprend également la prise en charge de différentes sources de données,
la gestion des erreurs, la mise en place de tests comparatifs et la préparation
d'un exécutable Windows autonome.

La problématique générale de ce travail peut ainsi être formulée de la manière
suivante :

\begin{quote}
\textit{Comment remplacer intégralement les traitements FME utilisés par Scan
par une solution Python autonome, tout en garantissant la fiabilité des
résultats et la compatibilité avec les différentes sources de données prises en
charge ?}
\end{quote}

Pour répondre à cette problématique, le travail a été organisé autour de
l'analyse de l'existant, du développement des nouveaux traitements, de leur
intégration dans l'architecture de Scan et de leur validation par comparaison
avec les résultats FME de référence.

Ce rapport est structuré en quatre chapitres. Le premier présente le contexte de
l'alternance, l'entreprise Isogeo, le produit Scan ainsi que les objectifs de la
mission. Le deuxième expose l'analyse de l'existant et les choix techniques
retenus pour concevoir la solution. Le troisième décrit les développements
réalisés, la préparation de l'exécutable et la démarche de validation. Enfin, le
quatrième présente les résultats obtenus, les apports de l'alternance, les
limites rencontrées et les perspectives du projet.


